% Authored appendix for the next draft. The source is deliberately independent
% of the main paper's private notation macros.
\section{The alarm construction and specified objectives}
\label{app:alarm}

This appendix supplies the literal reward definitions behind the comparison.
All returns below are expectations under the same geometric prior,
$\alpha(\infty)=1/2$ and $\alpha(k)=2^{-(k+2)}$, except where a fixed-world
quantity is explicitly used before averaging. The native comparison takes a
supremum over worlds and does not use this prior. Logarithmic rewards are reported in bits.
For nonnegative rewards, the complete return is the limit of the finite returns. A discounted return includes the startup decision with weight one.
These aggregation conventions are part of each objective.

\subsection{The interface and its exact native error}
Use the interface defined in Section~\ref{sec:alarm}, abbreviating INSPECT,
PLAY$_0$, and PLAY$_1$ by $\mathsf I,\mathsf P_0,\mathsf P_1$. We retain every
action and observation in the history. Let $b(\theta)$ be the inspection answer
in world $\theta$, and let $g(a)$ be the active display's digit after action $a$.
Thus $b(\theta)=\mathbf1\{\theta<\infty\}$ and
$g(a)=\mathbf1\{a=\mathsf P_1\}$.

Let $s=\pi(\mathsf I\mid\varnothing)$ and let $T_{\mathsf I}$ denote the
actual one-observation inspection experiment, with observation $(2b(\theta),0)$
(written $F$ in Section~\ref{sec:alarm}).
For a target experiment $T$, write $\delta(K_{\pi,t},T)$ for its deficiency
from the length-$t$ history acquired by $\pi$. The depth-$n$ native audit is the
supremum of this deficiency over the finite interventions of length at most
$n$ available on the interface.

\begin{theorem}[The common audit]\label{thm:next-alarm-audit}
For every randomized history policy $\pi$ and every $t\ge1$,
\[
 \delta(K_{\pi,t},T_{\mathsf I})=\frac{1-s}{2}.
\]
For every $t\ge n\ge1$, the depth-$n$ native audit has the same value.
Inspection simulates every policy's experiment at the same positive horizon.
Consequently $\pi$ is natively sufficient exactly when $s=1$; when $s<1$, an
inspection policy strictly dominates $\pi$ in the finitary process preorder.
\end{theorem}
\begin{proof}
Fix $t$ and compare $\infty$ with a finite world whose pulse comes after $t$.
Conditional on a play startup, their history laws agree, including the adaptive
action factors. Conditional on inspection, their root observations are respectively
$(0,0)$ and $(2,0)$. Since the initial action remains in the history, the total variation
distance of the two source rows is exactly $s$. The corresponding target rows
have distance one. If a decoder has error at most $\varepsilon$ in each world,
contraction and the triangle inequality give $1\le s+2\varepsilon$.

For the upper bound, extract the alarm coordinates in the source history. An inspected
history also gives the correct bit $b$. On a play history, guess $b$ fairly.
Choose a world consistent with the bit and the recorded alarm prefix, and
simulate the target with its own policy randomness and the fresh color bits
required on a play branch. If the
guessed bit is consistent and correct, all worlds still compatible with this
summary have identical target laws through depth $n\le t$. An impossible
summary can be filled arbitrarily. Thus inspection is exact, and the play
branch contributes error at most $1/2$. This proves the audit upper bound;
the startup target gives equality. The same decoder at $s=1$ proves
same-horizon sufficiency. A positive startup deficit at every source horizon
prevents reverse finitary domination.
\end{proof}

The lower bound concerns uniform simulation across worlds: the late-alarm world
may change with $t$, so the argument does not require one fixed finite alarm to
remain hidden forever. The upper bound likewise concerns simulated observations;
regenerating the display's random outputs from retained evidence does not make
the inspecting agent's physical display responsive again.

\subsection{Continuing pseudo-count bonuses at arbitrarily large horizons}
Let $O=\{0,1,2,3\}\times\{0,1\}$ be the raw observation alphabet, and
let $N_h(o)$ count occurrences of the whole pair $o\in O$ in a length-$t$
history $h$. A categorical probability model assigns one adjustable count
to each label. With initial count one on each of the eight labels, its probability
before and after one further occurrence of $o$ is
\[
 \rho_h(o)=\frac{N_h(o)+1}{t+8},\qquad
 \rho'_h(o)=\frac{N_h(o)+2}{t+9}.
\]
The density-update construction of \citet[Equations (1)--(2)]{bellemare2016unifying}
defines the effective count by
\[
 \widehat N_h(o)=\frac{\rho_h(o)(1-\rho'_h(o))}{\rho'_h(o)-\rho_h(o)}
               =N_h(o)+1.
\]
We reward $1/\sqrt{\widehat N_h(o)}$ before updating the model. This is the
inverse-square-root specialization with unit scale and the stated initialization;
it does not use their learned image model or numerical bonus offset. Let $S(n)$
be the cumulative bonus from $n$ occurrences of one label:
$S(n)=\sum_{j=1}^n j^{-1/2}$ and $S(0)=0$. Let $N_H(o)$ be the counts in a
length-$H$ history. Write $R_H$ for its cumulative reward and
$J_H^{\rm pseudo}(\pi)$ for its expected value under policy $\pi$. Then
\[
 R_H=\sum_{o\in O} S(N_{H}(o)),\qquad J_H^{\rm pseudo}(\pi)=\mathbb E R_H.
\]
Expectation includes the prior, observations, and policy randomization.

\begin{proposition}[Continuing count bonus]\label{prop:alarm-pseudocount}
On the fixed alarm interface, for every integer $H\ge64$, the objective
$J_H^{\rm pseudo}$ has an attained global maximum over all randomized history
policies. Every maximizer skips initial inspection. Every extension of a maximizing
length-$H$ policy has inspection-target deficiency $1/2$ at every positive collection
time and eventually, and is strictly finitarily dominated by sufficient inspection.
\end{proposition}
\begin{proof}
The elementary bounds
\[
 2\sqrt{n+1}-2\le S(n)\le2\sqrt n,\qquad S(a+b)\le S(a)+S(b)
\]
follow by telescoping square-root increments and decreasing summands.
On any inspection branch at most two observations differ from $(0,0)$: the startup
answer $(2,0)$ and the single possible alarm pulse $(0,1)$. All others are $(0,0)$,
regardless of subsequent actions. Hence every such branch satisfies
\[
 R_H\le S(H)+2\le2\sqrt H+2.
\]
By contrast, start with PLAY and then alternate PLAY$_0$ and PLAY$_1$.
Let $n_0$ and $n_1$ count the later PLAY$_0$ and PLAY$_1$ actions, respectively.
Among the $H-1$ later steps, $n_0=\lceil(H-1)/2\rceil$ and
$n_1=\lfloor(H-1)/2\rfloor$. There are at least $n_0-1$ observations in
$\{(0,0),(1,0)\}$ and at least $n_1-1$ in
$\{(2,0),(3,0)\}$, because at most one alarm coordinate is nonzero.
For $H\ge64$, subadditivity therefore gives, on every supported history,
\[
 R_H\ge S(n_0-1)+S(n_1-1)\ge4\sqrt{n_1}-4
 \ge2\sqrt{2H-4}-4>2\sqrt H+2.
\]
These bounds survive prior expectation. Finite alphabets and a finite reward horizon
make the objective a finite multilinear function of policy rows, so a deterministic
finite plan attains the global maximum. The displayed PLAY policy strictly outperforms
every inspection policy. Splitting any policy by the initial action therefore shows
that a maximizer assigns zero probability to inspection. Changing the maximizer's actions after
$H$ cannot change this conclusion. The native deficiency and strict domination follow
from Theorem~\ref{thm:next-alarm-audit}.
\end{proof}

Unlike finite returns at horizons $H\ge64$, the complete and long-run average
returns tie all policies. For every
length-$H$ observation sequence, subadditivity and the square-root bounds give
\[
 \sqrt H\le S(H)\le R_H\le2\sum_{o\in O}\sqrt{N_H(o)}\le4\sqrt{2H}.
\]
Thus $J_H^{\rm pseudo}(\pi)\to+\infty$ and
$J_H^{\rm pseudo}(\pi)/H\to0$ for every policy, including sufficient inspection.
Both limiting objectives tie all policies. This does not negate the proposition's
strict exclusion of inspection from every sufficiently long finite-horizon optimum.
The threshold $64$ is a convenient sufficient bound, not an optimal threshold.

\subsection{Information and Brier returns}
Here the world label $\theta$ is drawn from $\alpha$. Write $\beta_t$ for the
actual posterior over the countable world labels given the complete length-$t$
history. The information objective at horizon $t$
is $I(\theta;H_t)$; its complete value is $\sup_t I(\theta;H_t)$.
The Brier return sums
$\mathbb E\|\beta_{t+1}-\beta_t\|_2^2$. It uses posterior probabilities
of individual worlds, without replacing the geometric tail by one label.
The name comes from the category-summed quadratic forecast loss of
\citet[Equations (2)--(4)]{brier1950verification}, whose optimal expected loss
at a distribution $p$ is $1-\|p\|_2^2$. Expected reductions in this optimal loss
give the posterior-movement identity used here:
\[
 \sum_{j<t}\mathbb E\|\beta_{j+1}-\beta_j\|_2^2
 =\mathbb E\|\beta_t\|_2^2-\|\alpha\|_2^2.
\]
Here $H(\alpha)=2$ bits and $\|\alpha\|_2^2=1/3$.
Every policy's history yields its alarm prefix by deterministic extraction.
The information of those prefixes converges to two bits and their expected
squared posterior norm converges to one. Since the full history determines its
alarm prefix, these prefix values give lower bounds for each policy. Prior
entropy and the unit bound on the squared posterior norm supply the matching
upper bounds. Summing the successive expected score changes therefore
gives, for every policy,
\begin{equation}\label{eq:next-complete-posterior}
 \sup_t I(\theta;H_t)=2,\qquad
 \sum_{t\ge0}\mathbb E\|\beta_{t+1}-\beta_t\|_2^2=\frac23.
\end{equation}
Both complete objectives tie sufficient inspection with incomplete policies.

At each fixed positive finite horizon, however, inspection maximizes both
information and Brier returns, and equality holds exactly for
$s=1$. Theorem~\ref{thm:next-alarm-audit} supplies a same-horizon garbling
from inspection and rules out the reverse garbling when $s<1$. Strictness
of the actual full-support atomic information and Brier functionals gives
the strict preference. This positive result concerns the same countable
class; the change is from a complete return to a fixed finite horizon.

The entire infinite alarm stream determines the world: a pulse identifies its
finite delay, and an entirely silent stream identifies the never-alarm world.
The inspection answer can therefore be recovered from that infinite stream.
Every finite silent play record, however, is still compatible with a later alarm
and with no alarm. The uniform error in Theorem~\ref{thm:next-alarm-audit}
therefore persists at every finite collection time.

\subsection{Actual-world block empowerment with complete return}\label{app:alarm-empowerment}
A two-step probe just after startup sees two display--alarm pairs. In play mode,
each pair reveals the selected digit and adds an independent fair color, while
the alarm coordinates depend only on the world. The two controllable digits
give two bits of capacity; inspection instead makes both responses independent
of the actions and leaves zero capacity.

Fix a world $\theta$, a collected history $h$, and a positive integer $n$.
For an imposed action sequence $u\in A^n$ and raw observation sequence $y\in O^n$,
let $W_{\theta,h,n}(y\mid u)$ be the probability of $y$ when executing $u$
from $h$ under the actual response law. The probe does not adapt to intermediate
outputs. Its input distribution $\nu$ ranges over distributions on $A^n$.
Let $U$ and $Y$ denote the proposed action and resulting observation blocks.
Their joint law and the observation marginal are
\[
 P_{\nu,\theta,h,n}(u,y)=\nu(u)W_{\theta,h,n}(y\mid u),\qquad
 m_{\nu,\theta,h,n}(y)=\sum_u\nu(u)W_{\theta,h,n}(y\mid u).
\]
Define the physical channel capacity, in bits, by
\[
 C_{\theta,n}(h)=\max_\nu\sum_{u,y}\nu(u)W_{\theta,h,n}(y\mid u)
       \log_2\frac{W_{\theta,h,n}(y\mid u)}{m_{\nu,\theta,h,n}(y)}.
\]
Zero-probability terms contribute zero. Finite alphabets make the maximum
attainable and give $0\le C_{\theta,n}(h)\le n\log_2|A|$.
There is no posterior mixture inside this channel. The maximizing probe
law may differ between worlds, because capacity describes the physical
channel; the collecting policy remains one world-independent history rule.
The same world governs every step of a probe and the entire collected run.

At any nonempty history, every one of the next $n$ observations contains a display.
Consequently
\[
 C_{\theta,n}(h)=n\,\mathbf1\{h\text{ started in play mode}\}.
\]
Indeed, the observation block reveals the vector of proposed digit bits at
its $n$ steps. Conditional on that vector, the remaining randomness is
independent fair color; the alarm sequence is fixed by $\theta$ and the clock.
The channel factors through the digit vector, so its mutual information is at
most that vector's entropy, bounded by $n$ bits. Independent fair choices
of PLAY$_0$ and PLAY$_1$ at every step attain the bound. After inspection
all outputs are action-independent, hence capacity is zero. This also proves
that the full-block capacity at a nonempty history equals the sum of the
terminal-sensor capacities for probe lengths $1,\ldots,n$ on this particular
interface; that equality is not a general channel identity.

Write $s=\pi(\mathrm{INSPECT}\mid\varnothing)$ and
$b_n=\mathbb E_{\theta\sim\alpha}C_{\theta,n}(\varnothing)$.
The uniform bound on capacity makes $b_n$ finite and nonnegative.
For every randomized history policy and every $T\ge1$, the actual finite
expected return is
\begin{equation}\label{eq:actual-block-finite}
 R_{n,T}(\pi):=\mathbb E_{\alpha,\pi}\sum_{t=0}^{T-1}C_{\theta,n}(H_t)
 =b_n+n(1-s)(T-1).
\end{equation}
The expectation averages world-specific capacities after their evaluation,
and arbitrary policy and response randomness. Startup contributes the same
$b_n$ for every collecting policy. Each of the $T-1$ later histories contributes
$n$ in play mode and zero in inspection mode.

\begin{proposition}[Complete actual-world empowerment]\label{prop:actual-block-complete}
For each fixed positive probe length $n$, the nonnegative cumulative objective
$J_{{\rm emp},n}(\pi)=\sup_T R_{n,T}(\pi)$ satisfies
\[
 J_{{\rm emp},n}(\pi)=
 \begin{cases}b_n,&s=1,\\+\infty,&s<1.\end{cases}
\]
Its maximum is attained, and a policy maximizes it if and only if $s<1$.
Every maximizer is strictly finitarily dominated by inspection. Its deficiency
at every positive collection time and its eventual inspection deficiency are
$(1-s)/2$, and the set of optimal deficiencies is exactly $(0,1/2]$.
\end{proposition}
\begin{proof}
The coefficient of $1-s$ in~\eqref{eq:actual-block-finite} is unbounded for
every $n\ge1$, since it equals $n(T-1)$.
If $s=1$, every positive-horizon return is $b_n$. Otherwise the finite sums
diverge to infinity. A policy that always plays attains this infinite maximum.
The common audit gives $(1-s)/2$ and strict inspection domination whenever
$s<1$. Every value in $(0,1/2]$ is realized by randomizing at startup as follows: for a desired deficiency $d\in(0,1/2]$, choose $s=1-2d$. There is no positive uniform
lower bound on the deficiencies of all optimal policies.
\end{proof}

No time average, discount, or infinite-probe capacity is substituted into this
objective. At each fixed finite reward horizon $T\ge2$, the maximizers instead
have $s=0$; taking the complete sum introduces the additional infinite ties.
The full observation block is explicitly permitted by
\citet[Section~2.2, paragraph after Equation~(3)]{klyubin2005empowerment}.
The undiscounted cumulative return is the comparison chosen here; it is not
attributed to the source's learning or evolutionary algorithm.

The argument above is a written proof for the stated response law.
The formalization boundary, including the full-block capacity calculation and
the startup expectation, is recorded in Appendix~\ref{app:verification}.

\subsection{The native score on this interface}
Inspection is sufficient by Theorem~\ref{thm:next-alarm-audit}.
For the rich target family and strictly positive weights of
Theorem~\ref{thm:strict-finitary-objective}, the maximizing policies are therefore
exactly the sufficient policies:
\[
 \pi\in\operatorname*{arg\,max}J_w\quad\Longleftrightarrow\quad s=1.
\]
